% Generated draft: requires mathematical and source review.
\documentclass[11pt]{article}
\usepackage[margin=25mm]{geometry}
\usepackage{fontspec}
\setmainfont{Latin Modern Roman}
\usepackage{amsmath,amssymb,mathtools,unicode-math}
\setmathfont{Latin Modern Math}
\usepackage{enumitem,needspace,xurl,hyperref}
\setlength{\parindent}{0pt}
\setlength{\parskip}{5pt}
\setlength{\emergencystretch}{3em}
\newcommand{\N}{\mathbb{N}}
\newcommand{\Z}{\mathbb{Z}}
\newcommand{\Q}{\mathbb{Q}}
\newcommand{\R}{\mathbb{R}}
\newcommand{\C}{\mathbb{C}}
\newcommand{\F}{\mathbb{F}}
\newcommand{\A}{\mathbb{A}}
\newcommand{\PP}{\mathbb{P}}
\newcommand{\E}{\mathbb{E}}
\newcommand{\eps}{\varepsilon}
\newcommand{\dd}{\,\mathrm{d}}
\begin{document}
\section*{UnsolvedMath ID: 20002684}
\textbf{Source problem code:} catalogue entry
\subsection{Definitions and mathematical statement}
The missing display in the source is recovered from companion records as the real logarithmic determinant identity for a bounded commuting tuple:
\[
 \log\Delta\left(1-\sum_{i=1}^n\alpha_iT_i\right)
  =\int_{\C^n}\log\left|1-\sum_{i=1}^n\alpha_i\zeta_i\right|
                    \,\mu(\dd\zeta).
\]
The absolute-value bars are essential; a single branch of the complex logarithm cannot replace this real extended-valued expression globally. The question asks which functions $F$ on $\C^n$ have such an integral representation. The aid treats compactly supported probability measures and gives the following precise slice-compatibility criterion.

Let $F:\C^n\to[-\infty,\infty)$, require $F(0)=0$, and use the complex bilinear pairing $a\cdot\zeta=\sum_i a_i\zeta_i$. For each $a\in\C^n$, define on $\lambda\ne0$
\[
 U_a(\lambda)=\log|\lambda|+F(a/\lambda).
\]
Require that $U_a$ extend to a nontrivial subharmonic function on $\C$, with normalized behavior $U_a(\lambda)-\log|\lambda|\to0$ at infinity, whose Riesz measure
\[
 \nu_a=(2\pi)^{-1}\Delta_{\R^2}U_a
\]
is a compactly supported probability measure. These conditions identify $U_a$ with its logarithmic potential $\int\log|\lambda-w|\,\nu_a(\dd w)$. At $a=0$, $U_0=\log|\lambda|$ and $\nu_0=\delta_0$.

Impose two additional compatibility conditions. First, for one finite $R$, uniformly over all $a$,
\[
 \operatorname{supp}\nu_a\subseteq\{w:|w|\le R\|a\|_2\}.
\]
Second, the function
\[
 \Phi(a)=\int_\C e^{i\operatorname{Re}w}\,\nu_a(\dd w)
\]
is continuous and positive definite on the real additive group $\C^n$: for every finite choice $a_1,\ldots,a_k$ and $z_1,\ldots,z_k\in\C$,
\[
 \sum_{r,s=1}^k z_r\overline{z_s}\Phi(a_r-a_s)\ge0.
\]
The inequality means that this sum is real and nonnegative. The slice formula automatically gives the scaling consistency $\nu_{ca}=(w\mapsto cw)_*\nu_a$ for $c\in\C$, with the stated convention at $c=0$.

The aid reports that these conditions are necessary and sufficient for a unique probability measure supported in the closed Euclidean ball of radius $R$ to satisfy
\[
 F(\alpha)=\int_{\C^n}\log|1-\alpha\cdot\zeta|\,\mu(\dd\zeta)
 \quad(\alpha\in\C^n).
\]
Indeed Bochner's theorem reconstructs a joint measure from $\Phi$ using the real pairing $\operatorname{Re}(a\cdot\zeta)$; scaling consistency supplies all complex frequencies of each slice, so $(a\cdot\zeta)_*\mu=\nu_a$. The uniform all-direction support estimate gives the ball bound, not merely coordinatewise bounds. Conversely a compact joint measure produces exactly those slices and a continuous positive-definite characteristic function.

Every such measure is reported realizable as the joint spectral distribution of a commuting normal tuple in a $\mathrm{II}_1$ factor, by representing its bounded coordinate functions in an abelian subalgebra. For that tuple the determinant identity follows from normal functional calculus. This is a range criterion for the stated compact-probability transform; it is not a characterization of arbitrary multivariable subharmonic functions without joint slice compatibility.
\subsection{Short English statement}
Characterize logarithmic determinant transforms on $\C^n$. The aid reconstructs each one-dimensional slice by its Riesz measure and tests whether those slices come from one compact joint probability law through a continuous positive-definite characteristic function and a uniform directional support bound.
\subsection{Sources}
\begin{itemize}
\item American Institute of Mathematics, \emph{Free Analysis} problem list, physical page 6, subsection 0.9, ``Focus Group on Invariant Subspaces (day 4),'' \url{https://aimath.org/WWN/freeanalysis/freeanalysis.pdf}.
\item Local UnsolvedMath registry, permanent ID 20002684, source code \texttt{catalogue entry}, canonical location \texttt{aim-probability-notes.json}, \texttt{notes[125]}, record dated 14 August 2026. Its original question refers only to the omitted display, and the literature field is empty.
\item Companion permanent ID 20002682, \texttt{AIM-PROBABILITY-0124}, \texttt{notes[123]}, contains the display. Companion permanent ID 20002683, \texttt{AIM-PROBABILITY-0125}, \texttt{notes[124]}, identifies missing modulus signs by reference to Schultz's primary theorem and Charlesworth--Dykema--Sukochev--Zanin; full bibliographic entries are not supplied.
\item Supplied research attempt 1: partial result, problem status at run unknown, rigor labelled proved, novelty confidence low. Its slice--Bochner theorem is machine-generated research aid, not peer reviewed; independently verify before citation. Bochner's theorem is invoked without a particular bibliographic edition.
\end{itemize}
\end{document}
